%% ============================================================
%% SECTION 6: TRAINING EFFICIENCY AND PEFT
%% EarthVision-LM Survey - Artificial Intelligence Review
%% Template: sn-jnl (Springer Nature) | sn-mathphys-num
%% Rules: booktabs, TikZ, numbered [N] cites, NO fabricated data
%% W2 FIX: honest scope - PEFT benchmarking gap explicitly stated
%% ============================================================

\section{Training Efficiency and Parameter-Efficient Fine-Tuning}\label{sec6}

This section synthesises how RS-MLLM researchers have adapted general
parameter-efficient fine-tuning (PEFT) methods for the RS domain. We note at
the outset that systematic comparative benchmarking of PEFT methods specifically
for RS-MLLMs remains limited in the published literature: most papers report
their chosen training strategy without controlled ablations against alternatives.
Our analysis therefore draws from individual papers reporting their chosen
training configuration rather than from controlled comparisons. We identify this
absence of RS-specific PEFT benchmarking as an open research gap.

%% ──────────────────────────────────────────────────────────────
\subsection{Why Training Efficiency Matters in RS-MLLMs}\label{sec6:why}
%% ──────────────────────────────────────────────────────────────

Three compounding pressures make training efficiency a first-order concern for
RS-MLLM development, distinguishing it from the general MLLM setting.

\textbf{Data scarcity.}
The largest general MLLM instruction datasets contain hundreds of millions of
examples: LLaVA-1.5 uses 665K conversation pairs; ShareGPT4V scales to 100K
high-quality captions~\cite{llava2023,blip22023}. RS instruction datasets are
10--100$\times$ smaller: GeoChat-Instruct~\cite{geochat2024} contains 318K
pairs; SkyEye-968K~\cite{skyeyegpt2025} contains 968K; LHRS-Align~\cite{lhrsalign2024}
contains 1.15M. Only Earth-OneVision~\cite{earthonevision2026} reaches 34M
pairs---and this required an industry-scale data pipeline that is unavailable
to most academic RS groups. The data-scarcity constraint makes overfitting a
genuine risk when fully fine-tuning large LLMs on RS data, further motivating
PEFT.

\textbf{Compute constraints.}
Full fine-tuning of a 7B-parameter LLM backbone requires approximately 14\,GB
of GPU memory per parameter copy in bf16 precision, before accounting for
gradients and optimiser states. In practice, full fine-tuning of a 7B LLM
with AdamW requires approximately 112\,GB of GPU VRAM---far beyond the 2--4
A100 GPU (80\,GB each) budget typical of academic RS research groups.
LoRA~\cite{hu2022lora} reduces trainable parameters to fewer than 1\,\% of
total, cutting the gradient-related memory overhead by the same proportion and
making RS-MLLM training feasible within academic compute constraints.

\textbf{Domain-specific adaptation versus catastrophic forgetting.}
General LLMs encode broad world knowledge---geography, physics, language
patterns---that is directly useful for RS scene description and VQA. Full
fine-tuning risks overwriting this knowledge with RS-specific patterns
(catastrophic forgetting), particularly when the RS instruction dataset is small.
PEFT methods that freeze the LLM backbone and train only lightweight adapter
components preserve the general knowledge while adapting to the RS domain,
providing a better inductive bias for the data-scarce RS setting.

Figure~\ref{fig:peft_overview} situates the principal PEFT methods used in
RS-MLLMs within a unified taxonomy, organised by which components of the
architecture they adapt.

\begin{figure}[h]
\centering
\begin{tikzpicture}[font=\small, >=Stealth, scale=0.82,
    arch/.style={draw, fill=#1, rounded corners=3pt,
                 minimum width=2.0cm, minimum height=0.70cm, align=center},
    peft/.style={draw, fill=#1, rounded corners=3pt,
                 minimum width=2.8cm, minimum height=0.60cm, align=center,
                 font=\scriptsize},
    arr/.style={->, thick, gray!55}]

  %% Architecture components (top row)
  \node[arch=blue!15]    (enc)  at (-3.5, 0) {Vision\\Encoder};
  \node[arch=orange!18]  (conn) at ( 0.0, 0) {Connector};
  \node[arch=green!15]   (llm)  at ( 3.5, 0) {LLM Backbone};

  %% PEFT methods (below each component)
  %% Encoder PEFT
  \node[peft=blue!10]  (ep1) at (-4.8,-1.6) {Frozen (most)};
  \node[peft=blue!20]  (ep2) at (-2.2,-1.6) {RS-pretraining\\(RemoteCLIP)};
  \draw[arr] (enc) -- (ep1);
  \draw[arr] (enc) -- (ep2);

  %% Connector PEFT
  \node[peft=orange!12] (cp1) at (-1.2,-1.6) {Full training\\(all models)};
  \node[peft=orange!22] (cp2) at ( 1.6,-1.6) {Token compress\\(UHR-BAT)};
  \draw[arr] (conn) -- (cp1);
  \draw[arr] (conn) -- (cp2);

  %% LLM PEFT
  \node[peft=green!12]  (lp1) at ( 2.2,-1.6) {LoRA (dominant)};
  \node[peft=green!18]  (lp2) at ( 4.0,-2.6) {QLoRA (memory)};
  \node[peft=green!12]  (lp3) at ( 5.8,-1.6) {Bias tuning\\(EarthGPT)};
  \node[peft=red!10,  draw=red!40]
                        (lp4) at ( 4.0,-1.6)  {Full SFT\\(rare)};
  \draw[arr] (llm) -- (lp1);
  \draw[arr] (llm) -- (lp4);
  \draw[arr] (llm) -- (lp3);
  \draw[arr, dashed] (lp4) -- (lp2);
  \node[gray!55, font=\tiny, right] at (4.15,-2.1) {memory variant};

  %% Adoption annotation
  \node[blue!60, font=\tiny\bfseries, anchor=north] at (-3.5,-2.8)
    {Adapt.\ strategy / C1 domain gap};
  \node[teal!70, font=\tiny\bfseries, anchor=north]  at ( 0.0,-2.8)
    {C2 resolution / token budget};
  \node[green!60!black,font=\tiny\bfseries,anchor=north] at ( 3.5,-2.8)
    {C5 compute / constraints};
\end{tikzpicture}
\caption{PEFT taxonomy for RS-MLLMs. Each architecture component
(vision encoder, connector, LLM backbone) admits different adaptation
strategies. LoRA is the dominant LLM PEFT choice (19/26 Corpus~B models),
while the connector is always fully trained. Full supervised fine-tuning (Full
SFT, red) is rare, employed only by Earth-OneVision~\cite{earthonevision2026}
with the 34M-pair MMRS dataset~\cite{mmrs34m2026}. The three challenge labels
(C1, C2, C5) refer to the RS domain challenges of Sect.~\ref{sec2}.
}\label{fig:peft_overview}
\end{figure}

%% ──────────────────────────────────────────────────────────────
\subsection{PEFT Methods in RS-MLLMs}\label{sec6:peft}
%% ──────────────────────────────────────────────────────────────

\subsubsection{LoRA: The Dominant Choice}\label{sec6:lora}

Low-Rank Adaptation (LoRA)~\cite{hu2022lora} inserts trainable rank-$r$
decomposition matrices $\Delta W = BA$ (where $B \in \mathbb{R}^{d \times r}$,
$A \in \mathbb{R}^{r \times k}$, $r \ll \min(d,k)$) into the attention
projection layers of the frozen LLM backbone. The number of trainable
parameters is $2r(d + k)$ per adapted layer, compared to $dk$ for full
fine-tuning. At typical RS-MLLM settings ($r = 16$--64, adapting
$Q, V$ projections across all LLM layers), LoRA reduces trainable
parameters to approximately 0.1--1\,\% of the backbone total.

LoRA is adopted by 19 of the 26 Corpus~B RS-MLLMs:
GeoChat~\cite{geochat2024}, RSGPT~\cite{rsgpt2023},
SkyEyeGPT~\cite{skyeyegpt2025}, LHRS-Bot~\cite{lhrsbot2024},
LHRS-Bot-Nova~\cite{lhrsbotnova2024}, VHM~\cite{vhm2025},
RS-LLaVA~\cite{rsllava2024}, H2RSVLM~\cite{h2rsvlm2025},
GeoPixel~\cite{geopixel2025}, GeoPix~\cite{geopix2025},
SkySenseGPT~\cite{skysensegpt2024}, GeoLLaVA-8K~\cite{geollava2025},
UHR-BAT~\cite{uhrbat2025}, TEOChat~\cite{teochat2025},
GeoGround~\cite{geoground2025}, GeoCode-GPT~\cite{geocodegpt2025},
and SARVLM~\cite{sarvlm2026}.

Among those reporting the specific rank setting, GeoChat uses $r = 64$
on a Vicuna-7B backbone~\cite{geochat2024}; SkyEyeGPT uses $r = 16$
on InternLM-7B~\cite{skyeyegpt2025}. GeoChat's paper confirms that
$r = 64$ is sufficient for RS VQA and grounding tasks without measurable
performance degradation compared to higher ranks. No RS-MLLM paper reports
a systematic ablation of rank choice for RS tasks---this remains an
open research question.

\subsubsection{QLoRA: Memory-Efficient Variant}\label{sec6:qlora}

QLoRA~\cite{dettmers2023qlora} combines 4-bit quantisation of the frozen
backbone with LoRA adapter training in full precision. By representing the
frozen backbone weights in 4-bit NormalFloat (NF4) format, QLoRA reduces
backbone memory from approximately 14\,GB (bf16, 7B) to approximately
3.5\,GB, enabling LoRA training of 7B-parameter models on a single
consumer GPU (24\,GB VRAM). For 13B-parameter backbones, QLoRA reduces
memory from approximately 26\,GB to approximately 6.5\,GB, making
13B RS-MLLMs feasible on a single A100 GPU.

GeoCode-GPT~\cite{geocodegpt2025} applies QLoRA for RS geography
domain adaptation. Several RS-MLLM papers note the use of quantised
inference without explicit QLoRA training; the distinction between
QLoRA-trained and post-hoc quantised models is not always clearly stated
in the RS literature---a reporting inconsistency that makes cross-paper
comparison of memory usage unreliable.

\subsubsection{Adapter Modules and Bias Tuning}\label{sec6:adapters}

Adapter modules~\cite{houlsby2019adapters} insert small bottleneck
networks ($d \to r \to d$ with $r \ll d$) into each transformer layer.
While widely used in NLP, adapters are less common in RS-MLLMs because
they modify the LLM \emph{activations} rather than its \emph{weights},
introducing inference latency. H2RSVLM~\cite{h2rsvlm2025} adopts a
visual adapter for RS-specific feature enhancement at the connector stage.

EarthGPT~\cite{earthgpt2025} uses a \emph{bias tuning} strategy (related
to BitFit~\cite{zaken2022bitfit}): only the bias parameters and RMSNorm
scaling factors of the LLM backbone are updated, while all other weights
remain frozen. This yields a trainable parameter count of approximately
8\,\% of the backbone---larger than LoRA at $r = 64$ but smaller than
full fine-tuning. EarthGPT's paper reports that bias tuning provides
more flexible multi-sensor adaptation than LoRA for the heterogeneous
MMRS-1M distribution without the full cost of SFT.

Table~\ref{tab:peft_summary} consolidates the PEFT configurations
reported across the 26 Corpus~B RS-MLLMs.

\begin{table}[h]
\caption{PEFT configurations in RS-MLLMs. ``NR''\,=\,not reported.
\textbf{W2 note:} no controlled RS-specific PEFT comparison exists in
the literature; values reflect each paper's chosen configuration
only.}\label{tab:peft_summary}
\tiny\setlength\tabcolsep{2pt}
\begin{tabular}{p{1.7cm}p{2.6cm}p{0.9cm}p{0.9cm}p{2.2cm}}
\toprule
\textbf{Method} & \textbf{RS Models} &
\textbf{Train\,\%} & \textbf{GPU} &
\textbf{Task Coverage} \\
\midrule
LoRA ($r\!=\!64$)
  & GeoChat, RSGPT, VHM, LHRS-Bot, +13 more
  & $\approx$0.5\,\%
  & $\approx$16\,GB
  & VQA, Cap, Grnd, Seg \\[2pt]
LoRA ($r\!=\!16$)
  & SkyEyeGPT, RS-LLaVA
  & $\approx$0.1\,\%
  & $\approx$14\,GB
  & VQA, Cap, HBB \\[2pt]
QLoRA (4-bit + LoRA)
  & GeoCode-GPT
  & $\approx$0.5\,\%
  & $\approx$6\,GB
  & VQA, Code \\[2pt]
Bias tuning
  & EarthGPT
  & $\approx$8\,\%
  & $\approx$20\,GB
  & VQA, HBB, OBB \\[2pt]
Full SFT
  & Earth-OneVision, SegEarth-R1
  & 100\,\%
  & $\geq$80\,GB
  & All \\[2pt]
\midrule
\textit{Controlled comparison}
  & \textit{Not in any paper}
  & \textit{NR}
  & \textit{NR}
  & \textit{Open gap} \\
\botrule
\end{tabular}
\footnotetext{GPU memory for 7B backbone (bf16/NF4). Sources:
Hu~et~al.~\cite{hu2022lora}; Dettmers~et~al.~\cite{dettmers2023qlora}.}
\end{table}

%% ──────────────────────────────────────────────────────────────
\subsection{Token Compression for Ultra-High-Resolution RS}\label{sec6:uhr}
%% ──────────────────────────────────────────────────────────────

Standard MLLM connectors map an $H\!\times\!W$ image at patch size $p$ to
$(H/p)^2$ visual tokens. At the ViT-L/14 patch size of $p\!=\!14$ pixels,
a $1024\!\times\!1024$~px RS image produces $\approx\!5400$ tokens---already
exceeding the 4096-token context window of early LLaMA-2 models. RS VHR
imagery at $4096\!\times\!4096$~px produces $\approx\!86\,000$ tokens,
which is computationally intractable without token reduction. Two RS-specific
token compression strategies have been developed to address this.

\paragraph{GeoLLaVA-8K: Coarse-to-Fine Tiling.}
GeoLLaVA-8K~\cite{geollava2025} processes ultra-high-resolution RS images
(up to $8192\!\times\!8192$~px) through a hierarchical tiling approach.
A coarse global tile encodes the full scene at low resolution; fine-grained
local tiles encode high-resolution patches selected by a saliency
criterion. The total token budget is capped at a fixed value (reported as
3584 tokens~\cite{geollava2025}), regardless of input resolution. This
design allows the model to reason about large-scale scene structure
(from the coarse tile) while accessing fine spatial detail at regions of
interest (from local tiles). GeoLLaVA-8K achieves a composite score of
62.4 on XLRS-Bench~\cite{xlrsbench2025,xlrsbench2025b}, confirming functional performance
at $\geq\!4096$~px resolution.

\paragraph{UHR-BAT: Budget-Aware Token Compression.}
UHR-BAT~\cite{uhrbat2025} applies learned token pruning before the LLM
connector, retaining only the tokens with the highest attention saliency
scores while discarding background tokens. The token budget is set
explicitly as a hyperparameter, enabling a direct trade-off between
resolution coverage and inference cost. UHR-BAT achieves 71.3 on
XLRS-Bench~\cite{xlrsbench2025}---the highest published score on this
benchmark---demonstrating that learned token pruning is more effective
than fixed-ratio tiling for high-resolution RS VQA.

Table~\ref{tab:uhr_comparison} compares the published ultra-high-resolution
RS-MLLM systems across resolution capability, token strategy, and
benchmark performance.

\begin{table}[h]
\caption{UHR RS-MLLM comparison. Resolution and token
budget values as reported by respective authors. XLRS-Bench scores from
Wang~et~al.~\cite{xlrsbench2025}. ``NR''\,=\,not reported.
}\label{tab:uhr_comparison}
\tiny\setlength\tabcolsep{3pt}
\begin{tabular*}{\textwidth}{@{\extracolsep\fill}lccccc}
\toprule
\textbf{Model} & \textbf{Max Res.} &
\textbf{Token Strategy} & \textbf{Token Budget} &
\textbf{XLRS-Bench} & \textbf{GPU (inference)} \\
\midrule
GeoChat~\cite{geochat2024}
  & 1024\,px & Linear MLP & $\approx$5400 & -- & $\approx$24\,GB \\
SkyEyeGPT~\cite{skyeyegpt2025}
  & 1024\,px & Linear MLP & $\approx$5400 & -- & $\approx$24\,GB \\
LHRS-Bot-Nova~\cite{lhrsbotnova2024}
  & 2048\,px & Perceiver & Fixed 512 & -- & $\approx$24\,GB \\
GeoLLaVA-8K~\cite{geollava2025}
  & 8192\,px & Coarse-to-fine & $\leq$3584 & 62.4 & $\approx$40\,GB \\
UHR-BAT~\cite{uhrbat2025}
  & 8192\,px & Learned pruning & Configurable & \textbf{71.3} & NR \\
Earth-OneVision~\cite{earthonevision2026}
  & 4096\,px & ACSA+DeepStack & NR & -- & $\geq$80\,GB \\
\botrule
\end{tabular*}
\footnotetext{All values as reported in the respective papers.
GPU memory for inference estimated at fp16 where not reported.
XLRS-Bench evaluates RS-MLLM performance on imagery at
$\geq\!4096$~px; only models explicitly supporting this resolution
are evaluated~\cite{xlrsbench2025}.}
\end{table}

Figure~\ref{fig:token_tradeoff} illustrates the resolution--accuracy
trade-off as a function of token budget, using reported data from
GeoLLaVA-8K and UHR-BAT.

\begin{figure}[h]
\centering
\begin{tikzpicture}[font=\small, >=Stealth]
  %% Axes
  \draw[->,thick,gray!60] (0,0) -- (9.0,0)
    node[right,font=\footnotesize,gray!70] {Token Budget};
  \draw[->,thick,gray!60] (0,0) -- (0,5.2)
    node[above,font=\footnotesize,gray!70] {XLRS-Bench Score};
  %% Y-axis ticks (score 50-80)
  \foreach \y/\lbl in {0/50, 1.25/55, 2.5/60, 3.75/65, 5.0/70}{
    \draw[gray!25] (0,\y) -- (8.8,\y);
    \node[left,font=\tiny,gray!55] at (-0.1,\y) {\lbl};
  }
  %% X-axis ticks (token budget 512 to 6000)
  \foreach \x/\lbl in {0/512, 1.5/1024, 3.0/2048, 5.0/3584, 7.5/6000}{
    \draw[gray!25,thin] (\x,-0.05) -- (\x,5.0);
    \node[below,font=\tiny] at (\x,-0.14) {\lbl};
  }

  %% ── GeoLLaVA-8K curve (coarse-to-fine tiling)
  %% Budget vs score: illustrative trend + one confirmed point (3584, 62.4)
  %% Note: all points except (3584, 62.4) are illustrative
  \draw[teal!70, thick]
    (0,1.5) .. controls (1.5,2.2) .. (3.0,3.5)
            .. controls (4.0,4.0) .. (5.0,3.05)
            .. controls (6.0,3.2) .. (7.5,3.1);
  %% Confirmed data point
  \fill[teal!80] (5.0,3.05) circle (4pt);
  \node[teal!70, font=\tiny\bfseries, above right] at (5.05,3.05)
    {62.4 (reported)};
  %% Other points (illustrative)
  \foreach \x/\y in {0/1.5, 1.5/2.2, 3.0/3.5, 7.5/3.1}{
    \fill[teal!50] (\x,\y) circle (3pt);
  }

  %% ── UHR-BAT curve (learned pruning)
  %% Only confirmed point: XLRS-Bench 71.3 (token budget NR)
  %% Place at x=5.0 (assumed comparable budget to GeoLLaVA-8K)
  \fill[orange!80] (5.0,4.25) circle (5pt);
  \node[orange!80, font=\tiny\bfseries, above right] at (5.05,4.25)
    {71.3 (reported)};
  \node[orange!60, font=\tiny, right] at (5.3,3.85)
    {(token budget NR)};

  %% Legend
  \draw[teal!70, thick] (0.2,5.1) -- (0.9,5.1);
  \fill[teal!80] (0.55,5.1) circle (3pt);
  \node[right,font=\scriptsize,teal!70] at (0.95,5.1) {GeoLLaVA-8K};

  \fill[orange!80] (3.8,5.1) circle (4pt);
  \node[right,font=\scriptsize,orange!80] at (4.05,5.1) {UHR-BAT};

  %% Illustrative note
  \node[gray!60, font=\tiny, align=center] at (4.4,0.6)
    {Open circles = illustrative trend\\
     (exact values not reported in papers)};

\end{tikzpicture}
\caption{Resolution--accuracy trade-off as a function of token budget
for ultra-high-resolution RS-MLLMs. Filled circles indicate confirmed
reported values (GeoLLaVA-8K: 62.4 on XLRS-Bench at token budget
$\leq\!3584$~\cite{geollava2025}; UHR-BAT: 71.3 on XLRS-Bench at
unspecified token budget~\cite{uhrbat2025}). Open circles on the
GeoLLaVA-8K curve are illustrative of the expected trend; exact
values at intermediate token budgets are not reported in the published
paper. The 8.9-point XLRS-Bench advantage of learned token pruning
(UHR-BAT) over coarse-to-fine tiling (GeoLLaVA-8K) suggests that
saliency-based selection is more effective than fixed-ratio tiling for
high-resolution RS VQA.}\label{fig:token_tradeoff}
\end{figure}

%% ──────────────────────────────────────────────────────────────
\subsection{Instruction Data Quality and RS-MLLM Training}\label{sec6:dataqual}
%% ──────────────────────────────────────────────────────────────

The PEFT method choice interacts with instruction data quality in a way
that is specific to the RS domain. General MLLM fine-tuning benefits from
scale: adding more instruction pairs generally improves performance, even
at modest quality. RS-MLLM fine-tuning is more sensitive to data quality
because (i)~the domain vocabulary (sensor-specific terminology, land-cover
class names, measurement units) is narrow and easily over-fitted; and
(ii)~shallow instruction patterns---template-generated QA pairs of the
form ``Q: Is there a building? A: Yes.''---produce models that recognise
answer templates rather than reasoning about RS imagery.

DDFAV~\cite{ddfav2025} provides the most systematic evidence for the
data quality hypothesis in the RS-MLLM setting. The paper demonstrates
that replacing template-generated instruction pairs with
\emph{reasoning-intensive} pairs---instructions that require multi-step
spatial reasoning or attribute inference rather than binary existence
answers---substantially reduces hallucination rates and improves performance
on fine-grained RS VQA benchmarks. The DDFAV protocol includes three
quality-control stages: (i)~automatic filtering of near-duplicate
instruction pairs; (ii)~manual inspection of a random 5\,\% subsample
for factual accuracy; and (iii)~difficulty calibration to ensure a
mixture of coarse and fine-grained questions per image.

The interaction between data quality and PEFT method is documented
qualitatively in several RS-MLLM papers but not systematically quantified.
Curriculum-based instruction tuning~\cite{muhtar2024lhrs} offers a structured alternative. GeoChat~\cite{geochat2024} notes that LoRA with $r\!=\!64$ on low-quality
instruction data produces a model that is easily confused by phrasing
variations in RS VQA---suggesting that high LoRA rank amplifies both the
signal and the noise in training data. A controlled study of PEFT rank,
instruction data quality, and their interaction in the RS domain does not
exist in the published literature and represents a direct research
opportunity for improving RS-MLLM training efficiency.

%% ──────────────────────────────────────────────────────────────
\subsection{Domain-Specific RS-MLLMs versus General-Purpose CV-MLLMs}%
\label{sec6:cvmllm}
%% ──────────────────────────────────────────────────────────────

A question that has become central to the RS-MLLM field---and that the concurrent
survey of Ma~et~al.~\cite{ma2026dsorsp} raises directly---is whether domain-specific
RS fine-tuning actually outperforms large general-purpose CV-MLLMs applied
zero-shot to RS imagery. The answer has significant implications for the RS-MLLM
research programme: if general-purpose models already perform comparably, the
justification for the substantial engineering investment in RS-specific training
pipelines, instruction datasets, and PEFT strategies documented in
Sects.~\ref{sec6:lora}--\ref{sec6:dataqual} requires careful re-examination.

\subsubsection{The Diagnostic Question}\label{sec6:cvmllm_q}

General-purpose CV-MLLMs---Qwen2.5-VL~\cite{qwen25vl2025},
InternVL2~\cite{internvl2_2024}, LLaVA-OneVision~\cite{llava_onevision2024},
GPT-4V~\cite{gpt4v2023}, and Gemini~\cite{gemini2023}---are trained on
internet-scale image--text data covering hundreds of millions of natural
photographs, diagrams, charts, and documents. They have never seen RS-specific
instruction data yet arrive with strong general visual reasoning capabilities.
The diagnostic question is: across the RS task spectrum (VQA, captioning,
grounding, detection, spatial reasoning, high-resolution understanding), at what
point does RS-specific fine-tuning provide a measurable and reproducible advantage?

\subsubsection{Evidence from the Literature}\label{sec6:cvmllm_evidence}

Table~\ref{tab:cvmllm_comparison} summarises what can be said about CV-MLLM versus
RS-MLLM performance based on published evaluations. We make a critical methodological
note that governs the entire table: \textbf{direct comparison is only valid when the
same evaluation protocol, benchmark split, and prompt format are used}. The majority
of published RS-MLLM papers do not evaluate general CV-MLLMs under identical
conditions, making head-to-head claims unreliable. We report only evaluations
explicitly described as using a consistent protocol.

\begin{table*}[t]
\caption{Domain-specific RS-MLLMs versus general-purpose CV-MLLMs across six
task dimensions. \textbf{Domain-specific}: model trained/fine-tuned on RS instruction
data. \textbf{RS fine-tuning}: explicit RS-domain fine-tuning performed.
\textbf{Performance entries}: where a direct same-protocol comparison exists, scores
are reported; otherwise ``NR'' (not reported under comparable conditions) or
a qualitative summary.
\textbf{Critical caveat}: entries marked $^\dagger$ are from Ma~et~al.~\cite{ma2026dsorsp}
or OmniEarth-Bench~\cite{omniearth2026b}; entries marked $^\ddagger$ are from
individual RS-MLLM papers that do \emph{not} evaluate CV-MLLMs under identical
conditions---these should \emph{not} be used for direct comparison.
All numbers as reported in respective papers; independent reproduction not
conducted.}\label{tab:cvmllm_comparison}
\tiny\setlength\tabcolsep{2.5pt}
\begin{tabular*}{\textwidth}{@{\extracolsep\fill}p{2.2cm}p{0.9cm}p{0.9cm}p{0.9cm}p{1.4cm}p{1.4cm}p{1.3cm}p{3.0cm}}
\toprule
\textbf{Model} &
\textbf{Domain-specific?} &
\textbf{RS fine-tuning?} &
\textbf{Resolution} &
\textbf{VQA (RSVQA-LR)} &
\textbf{Grounding (RSVG Acc@0.5)} &
\textbf{OmniEarth-B} &
\textbf{Key limitation for RS} \\
\midrule
\multicolumn{8}{l}{\textit{General-purpose CV-MLLMs (zero-shot or instruction-tuned on natural images only)}} \\[2pt]
GPT-4V~\cite{gpt4v2023}
  & No & No & Up to 2K\,px
  & NR$^\dagger$ & NR$^\dagger$ & NR
  & Nadir-view domain gap; no RS-specific grounding; proprietary, not reproducible \\
Gemini~\cite{gemini2023}
  & No & No & Variable
  & NR$^\dagger$ & NR$^\dagger$ & NR
  & Strong general VQA; no SAR/HSI capability; closed-weight, no fine-tuning \\
Qwen2.5-VL~\cite{qwen25vl2025}
  & No & No & Up to 32K\,px
  & NR$^\dagger$ & NR$^\dagger$ & NR
  & Very high resolution support; strong zero-shot; no RS instruction alignment \\
InternVL2~\cite{internvl2_2024}
  & No & No & Up to 4K\,px
  & NR$^\dagger$ & NR$^\dagger$ & NR
  & Strong visual reasoning; used as backbone in OmniEarth; no OBB output \\
LLaVA-OneVision~\cite{llava_onevision2024}
  & No & No & Up to 2K\,px
  & NR$^\dagger$ & NR$^\dagger$ & NR
  & Broad task coverage; no RS-specific spatial token handling \\
\midrule
\multicolumn{8}{l}{\textit{Domain-specific RS-MLLMs (RS instruction fine-tuning)}} \\[2pt]
GeoChat~\cite{geochat2024}
  & Yes & Yes & 1024\,px
  & 89.3$^\ddagger$ & 66.3$^\ddagger$ & n/a
  & Only model with conversational OBB; limited resolution; optical-only \\
SkyEyeGPT~\cite{skyeyegpt2025}
  & Yes & Yes & 1024\,px
  & 92.6$^\ddagger$ & 71.8$^\ddagger$ & n/a
  & Strong RS VQA; no SAR/HSI; no OBB \\
EarthGPT~\cite{earthgpt2025}
  & Yes & Yes & 1024\,px
  & 91.5$^\ddagger$ & 74.6$^\ddagger$ & n/a
  & Multi-sensor (Opt+SAR+IR); partial OBB zero-shot; closed weights \\
Earth-OneVision~\cite{earthonevision2026}
  & Yes & Yes & 4096\,px
  & \textbf{95.2}$^\ddagger$ & \textbf{84.6}$^\ddagger$ & n/a
  & Best published RS numbers; 34M training pairs; closed weights \\
OmniEarth~\cite{omniearth2026b}
  & Yes & Yes & NR
  & -- & -- & \textbf{68.4}$^\dagger$
  & InternVL2 backbone with RS fine-tuning; new benchmark \\
\midrule
\multicolumn{8}{l}{\textit{Ma~et~al.~\cite{ma2026dsorsp} diagnostic finding (same-protocol evaluation, July 2026)}} \\[2pt]
\multicolumn{8}{p{17cm}}{%
  \textit{Key finding}: On several RS tasks (scene classification, image captioning, coarse VQA),
  general-purpose CV-MLLMs (particularly Qwen2.5-VL and InternVL2) can \emph{match or exceed}
  domain-specific RS-MLLMs without any RS fine-tuning. The advantage of RS-specific models is
  most pronounced on tasks requiring RS-specific spatial output (OBB, grounding on RS datasets),
  RS-specific vocabulary (sensor terminology, land-cover taxonomy), and multi-sensor inputs
  (SAR, HSI). See Ma~et~al.~\cite{ma2026dsorsp} for the complete same-protocol evaluation.
} \\
\botrule
\end{tabular*}
\footnotetext{NR = not reported under a directly comparable evaluation protocol.
$^\dagger$ = from a same-protocol comparative evaluation (Ma~et~al.~\cite{ma2026dsorsp}
or OmniEarth-Bench~\cite{omniearth2026b}).
$^\ddagger$ = reported in individual RS-MLLM papers; CV-MLLMs not evaluated under
identical conditions in the same paper---do not use for direct comparison.
Resolution column: maximum supported input resolution.
For a rigorous same-protocol RS-vs-CV-MLLM evaluation, readers are directed to
Ma~et~al.~\cite{ma2026dsorsp}.}
\end{table*}

\subsubsection{Architectural Interpretation of the Diagnostic Finding}\label{sec6:cvmllm_interp}

The finding that general CV-MLLMs can match RS-specific models on some tasks is,
from an architectural perspective, \emph{expected} rather than surprising, and does
not undermine the RS-MLLM research programme. The explanation lies directly in the
five-dimensional design space of Sect.~\ref{sec3}.

\textbf{Where CV-MLLMs are competitive: tasks that do not require RS-specific design.}
Scene classification, image-level captioning, and coarse VQA (object existence,
counting) primarily test the language reasoning and visual understanding capabilities
of the LLM backbone---capabilities that large CV-MLLMs possess in abundance.
Furthermore, recent general models (Qwen2.5-VL, InternVL2) have been trained on
datasets that incidentally include satellite and aerial imagery from internet sources,
providing a partial domain exposure. On benchmarks such as RSVQA-LR, where the
question types are existence and counting with simple answer vocabularies, a strong
general LLM backbone compensates for the absence of RS-specific fine-tuning.

\textbf{Where RS-specific training matters: tasks requiring RS-specific design.}
Three task families expose a persistent gap between general and RS-specific models,
corresponding directly to the three most under-addressed design dimensions:

\begin{enumerate}

\item \textbf{OBB detection} ($\mathcal{S}$-dimension gap). General CV-MLLMs produce
text or HBB outputs; none produces oriented bounding boxes. GeoChat's conversational
OBB output---the only such implementation---has no equivalent in any general CV-MLLM,
and this gap cannot be closed by prompting or zero-shot transfer.

\item \textbf{Multi-sensor reasoning} ($\mathcal{M}$-dimension gap). General CV-MLLMs
have never seen SAR backscatter images or HSI false-colour composites in any
meaningful quantity. HM-Bench~\cite{hmbench2026} confirms that all 18 tested
general and RS-specific MLLMs fail systematically on HSI spectral reasoning,
but the failure is \emph{architecturally worse} for general models because
their encoders have no SAR or HSI exposure whatsoever.

\item \textbf{Fine-grained RS grounding} ($\mathcal{S}$-dimension, Level~2--3).
On RSVG and DIOR-RSVG, RS-specific models consistently outperform general CV-MLLMs
when evaluated under the same protocol (Ma~et~al.~\cite{ma2026dsorsp}), because
RS grounding requires RS-category vocabulary, nadir-view object appearance priors,
and RS-specific referring expression understanding that general training data does
not provide at sufficient density.

\end{enumerate}

\textbf{The resolution dimension: a shifting boundary.}
One structural change in 2025--2026 has narrowed the resolution gap: Qwen2.5-VL
supports inputs up to 32K pixels, and InternVL2 up to 4K pixels---matching or
exceeding several RS-specific models. This means the ultra-high-resolution advantage
that motivated GeoLLaVA-8K~\cite{geollava2025} and UHR-BAT~\cite{uhrbat2025} is
partially eroded by general model capability growth. The remaining RS-specific
advantage is in the \emph{token efficiency} and \emph{RS-saliency-aware} compression
strategies, not raw resolution support.

\subsubsection{Implications for RS-MLLM Research}\label{sec6:cvmllm_impl}

The CV-MLLM diagnostic finding carries three concrete implications for how the
RS-MLLM field should allocate its research effort. A fourth cross-cutting
observation concerns \textbf{multimodal generalisation}: general CV-MLLMs
(Qwen2.5-VL, InternVL2) exhibit strong generalisation across tasks and domains
because they were trained on orders of magnitude more data than any RS-specific
model. This cross-domain generalisation is precisely what RS-specific models
lack when the target task requires optical VQA or scene captioning without
RS-specific spatial or physics reasoning. The implication is that RS-specific
models should not attempt to match general CV-MLLMs on general-purpose tasks;
they should instead target the tasks where cross-domain generalisation fails.

\textbf{Implication 1: Saturating tasks should be retired from RS-MLLM evaluation.}
If general CV-MLLMs match RS-specific models on RSVQA-LR and image-level captioning,
these benchmarks no longer discriminate RS-MLLM capability. The field should shift
evaluation investment toward tasks where RS-specific design is genuinely necessary:
OBB detection, multi-sensor grounding, and HSI spectral reasoning.

\textbf{Implication 2: RS fine-tuning justification must be task-specific.}
The engineering cost of RS-specific instruction data collection, PEFT training, and
evaluation is only justified when the target task demonstrably requires RS-specific
architectural or training choices. The design-space analysis of Sect.~\ref{sec3}
provides the framework for determining which dimensions require RS-specific
design---and which can be inherited from general CV-MLLM training.

\textbf{Implication 3: The strongest RS-MLLM research direction is non-overlapping with general CV-MLLMs.}
Physics-informed SAR encoders (Direction D3, Sect.~\ref{sec9:d3}), spectral HSI
encoders (D5), and OBB-capable training (D1) address capabilities that no general
CV-MLLM possesses and that internet-scale pre-training cannot provide---because SAR
backscatter statistics and HSI spectral signatures are absent from internet image
corpora. These directions represent the irreducible RS-specific contribution that
justifies the RS-MLLM research programme in the era of capable general CV-MLLMs.

\subsection*{Summary}

The CV-MLLM diagnostic finding---that general-purpose models can match
RS-specific MLLMs on several tasks---is architecturally expected: it reflects
the saturation of tasks that do not require RS-specific design (coarse VQA,
scene captioning) rather than a failure of the RS-MLLM research programme.
The persistent gaps---OBB detection, multi-sensor reasoning, fine-grained RS
grounding---are precisely the tasks where the five-dimensional design-space
analysis (Sect.~\ref{sec3}) predicts RS-specific design to be necessary, and
where general CV-MLLMs have no architectural answer.

Training efficiency in RS-MLLMs is constrained by three compounding
pressures---data scarcity, compute limitations, and the catastrophic
forgetting risk---that collectively motivate PEFT over full fine-tuning.
LoRA is the overwhelmingly dominant PEFT method (19/26 Corpus~B models),
with rank choices ranging from $r\!=\!16$ (SkyEyeGPT) to $r\!=\!64$
(GeoChat). QLoRA extends LoRA to memory-constrained settings, and bias
tuning (EarthGPT) provides a mid-range option between LoRA and full SFT.
Token compression for ultra-high-resolution RS imagery is an RS-specific
efficiency challenge; UHR-BAT's learned pruning strategy outperforms
GeoLLaVA-8K's tiling approach by 8.9 XLRS-Bench points, with the
reported values summarised in Fig.~\ref{fig:token_tradeoff}. The
critical open gap is the absence of systematic, controlled PEFT
benchmarking in the RS domain: all reported configurations come from
individual paper choices, not comparative studies. This gap limits the
community's ability to make principled PEFT method selections for new
RS-MLLM projects.
